\documentclass[12pt]{article}
\usepackage[utf8]{inputenc}
\usepackage[T1]{fontenc}
\usepackage{hyperref}
\title{Yerong Language Resources for K-12 Curriculum Development}
\author{Compiled by Claude Sonnet 5.5 (claude-sonnet-5-5) for Andr\'as Kornai}
\date{}
\begin{document}
\maketitle

\section{Introduction}
Yerong (ISO 639-3 yrn) is a Kra-Dai language of south-west Guangxi, China. Its speakers live in Longhe and Pohe townships of Napo County. Wikipedia treats it as the southern branch of Buyang, which is spoken in Guangnan and Funing counties of Yunnan. Speaker estimates in the sources range from about 300--400 (1990) to 660 (2010). The Chinese state counts the Yerong as part of the Yao nationality. Searches found no school or curriculum materials.

\section{Primary Educational Resources}
None were found.

\section{Linguistic and Reference Materials}
\begin{itemize}
\item Wikipedia, Buyang language (covers Yerong as a branch): \url{https://en.wikipedia.org/wiki/Buyang_language}
\end{itemize}
This is a descriptive overview. I did not check it for further references.

\section{Religious and Community Resources}
\begin{itemize}
\item Joshua Project profile, Yerong in China (a Christian mission-oriented page with population figures): \url{https://joshuaproject.net/people_groups/16028}
\item Asia Harvest people-group profile, Yerong, Guangxi: \url{https://people-groups.asiaharvest.org/people-group-profiles/guangxi/yerong.pdf}
\end{itemize}
Both are mission-organisation profiles. Their figures should be checked against academic sources.

\section{Digital and Audio Resources}
None were found.

\section{Curriculum Development Context}
There is no known written standard or school use. Any curriculum would need new documentation of the language and the involvement of the community in Napo County. Work on the better-documented Buyang varieties in Yunnan may help.
\end{document}